\documentclass[../../../include/open-logic-section]{subfiles}

\begin{document}

\olfileid[fr]{sth}{story}{extensionality}
\olsection{Extensionnalité}

La première chose à dire est que ce sont leurs !!{element}s qui
permettent de distinguer les ensembles. Plus précisément :

\begin{axiom}[Extensionnalité]
Si les ensembles $A$ et $B$ ont les mêmes !!{element}s, alors $A$ et $B$
sont le même ensemble.
\[
  \lforall[A][\lforall[B][(\lforall[x][(x \in A \liff x \in B)] \lif
  \eq[A][B])]]
\]
\end{axiom}

Nous avons admis ce principe tout au long de la \olref[sfr][][]{part}.
Il mérite pourtant d'être rappelé. L'axiome d'extensionnalité exprime
l'idée fondamentale qu'un ensemble est déterminé par ses !!{element}s.
(On peut ainsi opposer les ensembles aux \emph{concepts} : les mêmes
objets peuvent relever de nombreux concepts distincts.)

Pourquoi adopter ce principe ? On peut soutenir que rejeter
l'extensionnalité reviendrait à renoncer à toute théorie qu'on pourrait
encore appeler \emph{théorie des ensembles}. La théorie des ensembles
n'est rien d'autre que la théorie des collections extensionnelles.

Le véritable défi, dans la \olref[sth][][]{part}, consiste toutefois à
formuler des principes qui nous disent \emph{quels ensembles existent}.
Or la seule réponse vraiment «~évidente~» à cette question est une
réponse dont on peut démontrer qu'elle est fausse.

\end{document}
